%! Skills Focus: Selecting an Appropriate Inference Procedure for Categorical Data — Unit 8, Lesson 8.7 — Lesson Plan
\documentclass[10pt]{article}
\usepackage{apstats-article}
\usepackage{apstats-boxes}
\usepackage{graphicx}
\graphicspath{{images/}}
\newcommand{\TallMath}[1]{$\displaystyle #1\rule[-1.4em]{0pt}{3.2em}$}
\newcommand{\UnitNumberName}{Unit 8: Inference for Categorical Data: Chi-Square \quad}
\newcommand{\LessonNumberName}{Lesson 8.7: Skills Focus --- Selecting an Appropriate Inference Procedure for Categorical Data}

\begin{document}
\begin{center}
    {\Large\bfseries \CourseName: \SchoolYear \\
    {\small\bfseries \UnitNumberName \LessonNumberName}}
\end{center}

\begin{tcolorbox}[colback=sky, colframe=navy, boxrule=0.9pt, arc=2mm,
    left=2mm, right=2mm, top=1.5mm, bottom=1.5mm]
    \textbf{Primary Objective:} Students will identify the correct inference procedure for any
    categorical data scenario by distinguishing among the five categorical inference procedures
    (z-tests/intervals for proportions, chi-square goodness-of-fit, chi-square homogeneity, and
    chi-square independence), and will explain why each procedure is or is not appropriate.
    (Big Idea: VAR, UNC)
\end{tcolorbox}

\renewcommand{\arraystretch}{1.6}
\begin{skillbox}[Priority Ideas \& Skills]{goldbox}
\begin{minipage}[t]{0.48\linewidth}
    \textbf{AP Skill 1 --- Selecting Statistical Methods}
    \begin{itemize}
        \item Identify the appropriate categorical inference procedure for a given scenario
              by recognizing the parameter structure: one proportion, difference of proportions,
              or a categorical distribution (Skill 1.E; VAR-6, UNC-4).
        \item Distinguish between a z-procedure (one or two proportions) and a chi-square
              procedure (three or more categories, or a two-way table) (Skill 1.D).
        \item Recognize whether a scenario calls for a confidence interval vs.\
              a significance test (Skill 1.D).
    \end{itemize}
\end{minipage}
\hfill
\begin{minipage}[t]{0.48\linewidth}
    \textbf{Key Understandings}
    \begin{itemize}
        \item Two-category binary questions $\to$ z-procedure. Three or more categories, or a
              two-way table $\to$ chi-square.
        \item Chi-square GOF: one variable, comparing an observed distribution to a claimed one.
        \item Chi-square homogeneity: one categorical variable, compare across two or more groups.
        \item Chi-square independence: two categorical variables recorded on one sample.
        \item Homogeneity and independence use the same test statistic and mechanics, but
              hypotheses and conclusions differ based on study design.
    \end{itemize}
\end{minipage}
\end{skillbox}

\begin{skillbox}[{Vocabulary, Concepts, \& Theorems}]{greenbox}
\textbf{Categorical Inference Procedure Selection Table}

\medskip
\renewcommand{\arraystretch}{1.8}
\begin{tabularx}{\linewidth}{@{} p{0.17\linewidth} p{0.22\linewidth} p{0.14\linewidth} X @{}}
  \rowcolor{sky}
  \textbf{Parameter / Structure} & \textbf{Trigger Words} & \textbf{Goal} & \textbf{Procedure Name} \\
  \hline
  One proportion $p$ & one group, binary outcome & CI & One-sample $z$-interval for $p$ \\
  One proportion $p$ & one group, binary outcome & Test & One-sample $z$-test for $p$ \\
  $p_1 - p_2$ & two groups, binary outcome & CI & Two-sample $z$-interval for $p_1-p_2$ \\
  $p_1 - p_2$ & two groups, binary outcome & Test & Two-sample $z$-test for $p_1-p_2$ \\
  Distribution over $\ge 3$ categories & one sample, claimed distribution & Test & Chi-square goodness-of-fit test \\
  One categorical variable & two or more independent groups & Test & Chi-square test for homogeneity \\
  Two categorical variables & one sample, two-way table & Test & Chi-square test for independence \\
\end{tabularx}
\end{skillbox}

\begin{skillbox}[Activate Prior Knowledge \& Spiral Review]{sky}
\begin{minipage}[t]{0.55\linewidth}\vspace{0pt}
    \textbf{Spiral Review (warm-up)}
    \begin{itemize}
        \item Names and key conditions for all five categorical inference procedures
              (Units 6 and 8)
        \item Recognizing parameter type from scenario language
        \item Distinguishing binary outcomes from multi-category outcomes
    \end{itemize}
\end{minipage}
\hfill
\begin{minipage}[t]{0.38\linewidth}\vspace{0pt}\centering
    % Authored warm-up compiles to target/ — no source PDF to embed here.
\end{minipage}
\end{skillbox}

\begin{skillbox}[Hook]{sky}
\begin{minipage}[t]{0.55\linewidth}
    \textbf{Which Procedure? (3--4 min)}\\
    Display two scenarios side by side:\\[4pt]
    A: ``A researcher surveys 200 students and records their favorite season (Fall, Winter,
    Spring, Summer). Is this consistent with equal preference for all seasons?''\\[4pt]
    B: ``A researcher surveys 200 students and records whether each prefers warmer months
    ($\ge$ Spring/Summer) or cooler months (Fall/Winter). Is the proportion who prefer
    warmer months equal to 0.5?''\\[4pt]
    Ask: ``Both involve seasons and preferences. Why are these different procedures?''
    Facilitate a 60-second Think-Pair-Share.
\end{minipage}
\hfill
\begin{minipage}[t]{0.40\linewidth}
    \textbf{The Key Diagnostic Questions}
    \begin{enumerate}
        \item \textbf{How many categories?}\\
              Two (binary) $\to$ z-procedure\\
              Three or more $\to$ chi-square
        \item \textbf{How many variables / groups?}\\
              One variable, one sample: GOF or z-interval/test\\
              One variable, two+ groups: homogeneity\\
              Two variables, one sample: independence
        \item \textbf{Estimate or test?}\\
              ``What is the range?'' $\to$ CI (z only)\\
              ``Is there evidence for a claim?'' $\to$ test
    \end{enumerate}
\end{minipage}
\end{skillbox}

\begin{skillbox}[Lesson: The Decision Framework]{sky}
\begin{minipage}[t]{0.48\linewidth}
    \textbf{Step-by-Step Procedure Selection}
    \begin{enumerate}
        \item \textbf{Count the categories.}\\
              Binary (success/failure, yes/no): z-procedure.\\
              Three or more: chi-square.
        \item \textbf{Identify the design.}\\
              Single sample + claimed distribution $\to$ GOF.\\
              Two+ independent groups, one categorical variable $\to$ homogeneity.\\
              Two categorical variables, one sample $\to$ independence.
        \item \textbf{Decide: CI or test?}\\
              CI applies only to z (proportion) procedures.\\
              Chi-square procedures are always significance tests.
        \item \textbf{Name the procedure} exactly.
    \end{enumerate}
\end{minipage}
\hfill
\begin{minipage}[t]{0.48\linewidth}
    \textbf{Homogeneity vs.\ Independence --- How to Tell}
    \begin{itemize}
        \item \textbf{Homogeneity:} Data come from \emph{separate} random samples (one per
              group). The question is whether the distribution of a categorical variable is the
              same across groups.
              ``Were freshmen and seniors randomly selected separately, then asked about
              preferred communication style?'' $\to$ homogeneity.
        \item \textbf{Independence:} Data come from \emph{one} random sample; two categorical
              variables are recorded for each subject.
              ``Was one random sample of students taken, and both grade level and communication
              preference recorded for each?'' $\to$ independence.
        \item Same $\chi^2$ statistic and table; different hypotheses and conclusions.
    \end{itemize}
\end{minipage}
\end{skillbox}

\begin{skillbox}[Explicit Instruction: Hypotheses and Conditions by Procedure]{sky}
\renewcommand{\arraystretch}{1.9}
\begin{tabularx}{\linewidth}{@{} p{0.22\linewidth} X X @{}}
  \rowcolor{sky}
  \textbf{Procedure} & \textbf{$H_0$ form} & \textbf{Key Condition (beyond Random/10\%)} \\
  \hline
  One-sample $z$-test & $H_0: p = p_0$ & Large Counts: $np_0 \ge 10$ and $n(1-p_0) \ge 10$ \\
  Two-sample $z$-test & $H_0: p_1 = p_2$ & Large Counts: all four $n\hat p$ and $n(1-\hat p) \ge 10$ \\
  Chi-square GOF & $H_0$: the distribution of [variable] is [claimed distribution] & Expected counts $\ge 5$ in each cell \\
  Chi-square homogeneity & $H_0$: the distribution of [variable] is the same across all groups & Expected counts $\ge 5$ in each cell \\
  Chi-square independence & $H_0$: there is no association between [var 1] and [var 2] & Expected counts $\ge 5$ in each cell \\
\end{tabularx}
\end{skillbox}

\begin{skillbox}[Active Monitoring]{redbox}
\begin{minipage}[t]{0.48\linewidth}
    \textbf{Circulate and Check}
    \begin{itemize}
        \item Students name the correct procedure and can articulate why (binary vs.\ multi-category;
              one group vs.\ two+ groups vs.\ one sample with two variables).
        \item Hypotheses use the correct parameter language for each procedure.
        \item Students do not confuse homogeneity with independence.
        \item Conditions are checked with numbers, not merely stated as ``met.''
    \end{itemize}
\end{minipage}
\hfill
\begin{minipage}[t]{0.48\linewidth}
    \textbf{Cold-Call Prompts}
    \begin{itemize}
        \item ``There are four categories --- which class of procedures does that trigger?''
        \item ``The researcher took one sample and recorded two variables --- homogeneity or
              independence?''
        \item ``Why can't we use chi-square to produce a confidence interval?''
        \item ``State $H_0$ for a chi-square test of independence.''
    \end{itemize}
\end{minipage}
\end{skillbox}

\begin{skillbox}[Group Work \& Differentiation]{redbox}
\begin{multicols}{3}
    \textbf{Tier R --- Remediate}
    \begin{itemize}
        \item Provide the selection table as a reference.
        \item Activity: identify procedure and state the form of $H_0$ only.
        \item Focus on the binary vs.\ multi-category diagnostic.
    \end{itemize}
    \columnbreak
    \textbf{Tier A --- Approaching}
    \begin{itemize}
        \item Complete all six activity scenarios (identify, state hypotheses, check conditions).
        \item Distinguish homogeneity from independence in their own words.
    \end{itemize}
    \columnbreak
    \textbf{Tier E --- Extension}
    \begin{itemize}
        \item Complete all scenarios plus critique a provided student response.
        \item Explain when a two-sample $z$-test is preferred over chi-square homogeneity
              even when both are technically valid.
    \end{itemize}
\end{multicols}
\end{skillbox}

\begin{skillbox}[Individual Work \& Assessment]{redbox}
\begin{minipage}[t]{0.48\linewidth}
    \textbf{Exit Ticket (5 min)}\\
    Two to three brief scenarios: students name the correct procedure and state $H_0$.
    \begin{enumerate}
        \item A researcher surveys a random sample of 300 college students, recording gender
              (male/female) and whether each owns a car (yes/no). Test for an association.
        \item A candy company claims its bags contain 30\% red, 20\% orange, 20\% yellow,
              20\% green, and 10\% blue pieces. A student counts colors in one bag to test
              this claim.
    \end{enumerate}
\end{minipage}
\hfill
\begin{minipage}[t]{0.48\linewidth}
    \textbf{AP-Style Check}\\
    \textit{A political scientist randomly selects 150 voters from each of three regions and asks
    each voter's party affiliation (Democrat, Republican, Independent). She wants to test
    whether party affiliation distribution is the same in all three regions. Which procedure?}
    \begin{enumerate}[label=(\Alph*)]
        \item One-sample $z$-test for $p$
        \item Two-sample $z$-test for $p_1-p_2$
        \item Chi-square goodness-of-fit test
        \item Chi-square test for homogeneity \textcolor{keyred}{\textbf{$\leftarrow$ correct}}
    \end{enumerate}
\end{minipage}
\end{skillbox}

\begin{skillbox}[Reinforcement \& Extension]{goldbox}
\begin{minipage}[t]{0.48\linewidth}
    \textbf{Homework Overview}
    \begin{itemize}
        \item Five AP-style scenario identification problems spanning all five categorical
              inference procedures; at least one requires a full four-step set-up (state + plan).
        \item One problem asks students to compare two procedures that could both apply and
              explain which is preferred and why.
        \item Extension: why chi-square procedures do not produce confidence intervals.
    \end{itemize}
\end{minipage}
\hfill
\begin{minipage}[t]{0.48\linewidth}
    \textbf{Spiral Review}
    \begin{itemize}
        \item All five categorical procedure names and triggering data structures
        \item Condition names for chi-square vs.\ z-procedures
        \item Hypothesis language for each procedure type
    \end{itemize}
    \textbf{Preview}\\
    Unit 8 concludes with the Sample Test --- students will apply all categorical inference
    procedures under timed, AP-exam conditions.
\end{minipage}
\end{skillbox}

\end{document}
